\documentclass[10pt,twocolumn]{article}
\usepackage[letterpaper,top=0.64in,bottom=0.68in,left=0.66in,right=0.66in,columnsep=0.24in]{geometry}
\usepackage[T1]{fontenc}\usepackage{lmodern}\usepackage{microtype}\usepackage{graphicx}\usepackage{booktabs}
\usepackage{array}\usepackage{caption}\usepackage{xcolor}\usepackage{enumitem}
\usepackage{balance}
\usepackage[hidelinks]{hyperref}\usepackage{url}
\definecolor{sprited}{HTML}{6337C7}\definecolor{teal}{HTML}{008C95}\definecolor{pale}{HTML}{F3F0FA}\definecolor{muted}{HTML}{555555}
\hypersetup{colorlinks=true,linkcolor=sprited,citecolor=sprited,urlcolor=sprited,
pdftitle={Co-Generating the Rig: Structure Tokens Improve Block-Autoregressive Video Diffusion},
pdfauthor={Jin Hyuk Cho}}
\captionsetup{font=small,labelfont=bf,skip=4pt}\setlist{nosep,leftmargin=*}
\setlength{\parindent}{1em}\setlength{\parskip}{0pt}\setlength{\textfloatsep}{8pt plus 2pt minus 2pt}
\setlength{\floatsep}{7pt plus 2pt minus 2pt}\setlength{\intextsep}{7pt plus 2pt minus 2pt}\renewcommand{\arraystretch}{1.08}

% TODO-block helper: unresolved evidence slots are typeset loudly so a draft
% compile can never silently pass for a finished paper.
\newcommand{\todoblock}[1]{\begin{center}\fcolorbox{sprited}{pale}{\parbox{.94\columnwidth}{\small\textbf{TODO:} #1}}\end{center}}

\title{\vspace{-1.5em}\textbf{Co-Generating the Rig: Structure Tokens Improve\\Block-Autoregressive Video Diffusion}}
\author{Jin Hyuk Cho\\Sprited\\\small\href{mailto:jin@sprited.ai}{jin@sprited.ai} $\cdot$
\href{https://orcid.org/0009-0001-1896-8242}{ORCID 0009-0001-1896-8242}}
\date{}

\begin{document}\raggedbottom\maketitle\vspace{-1.25em}

\begin{abstract}
% TODO(final pass): tighten once SRE and the sweep verdict land.
Block-autoregressive diffusion is the practical recipe for streaming video generation, and recent large-scale systems
report that jointly learning auxiliary modalities improves generation. We study this question causally, in the
setting where it matters most — teacher-forced block-autoregressive continuation — on a fully traced synthetic
testbed~\cite{cho2026dsf} whose ground-truth rig is exact. Our headline result: making the model \emph{co-generate an
explicit rig} (one 2D-skeleton token per latent frame, denoised jointly with the pixels under a shared flow) improves
pixel-space next-block prediction at matched parameters, data, and budget, on both teacher-forced and free-running
rollouts, at two training budgets. % TODO: exact divergence numbers once the sweep verdict freezes the table.
We measure this with an AR-native evaluation — next-block divergence against the data's actual continuation, with a
same-prompt real-pair floor — and, via a rig-space regressor (SRE), score rig-free and rig-generating models in the
same structural space. % TODO: SRE results.
Two supporting studies complete the report. First, we diagnose and repair a training pathology of this factorization:
a \emph{static-copy shortcut} where motion collapses while structure converges; history-side remedies fail while
signal-reallocation remedies fix it, and a matched full-clip control attributes the collapse to the factorization
itself. Second, we chart how tightly the rig should be coupled to the pixels: a differentiable soft-render coupling
binds rig to pixels an order of magnitude faster than co-generation alone, but its weight trades directly against
pixel quality. % TODO: sweep verdict sentence.
\end{abstract}

\section{Introduction}
\todoblock{Write after evidence sections stabilize. Frame: (1) block-AR diffusion is the streaming recipe
(MAGI-1, Diffusion/Self Forcing); (2) scale systems observe that multi-modal joint training helps generation
(UnityVideo, CVPR 2026: skeleton/depth/flow co-training lowers RGB loss at 10B/1.3M pairs) but full-clip settings,
uncontrolled scale, and shallow attribution leave the causal question open — and nothing measures continuation;
(3) we answer it in a controlled block-AR testbed: explicit rig co-generation, matched controls, pre-declared
protocols, AR-native next-block divergence; (4) contributions: causal evidence that rig co-generation improves
next-block prediction; the divergence metric + SRE rig-space evaluation; the static-copy diagnosis/repair; the
coupling-weight study.}

% --- moved from testbed report (Pixel P1): full-clip vs block-AR motivation essay ---
\paragraph{Why full-clip bidirectionality stops being enough.} Bidirectional temporal attention is useful inside the fixed 50-frame sample of K1, the testbed's pixel-space full-clip video DiT~\cite{cho2026dsf}, because every frame can coordinate with both earlier and later frames. The limitation
appears first in training cost. Every update must materialize and backpropagate through all 50 noisy output frames;
increasing the trained horizon grows activation memory across the network and the temporal-attention work
quadratically. In our uncompressed pixel setting, a longer horizon therefore quickly reduces batch size and makes each
optimization step substantially slower. The same interface also does not support native sliding-window continuation:
K1 jointly noises and predicts the entire window, rather than freezing an overlapping clean prefix and applying noise
and loss only to newly generated frames. For motion-dominated video, the latter is also an easier teacher-forced
learning problem: clean history already constrains identity, position, pose, velocity, and action phase, leaving the
model to denoise a smaller and less ambiguous future block rather than invent the full trajectory jointly. This
training advantage does not guarantee easier rollout, because generated history introduces exposure bias and error
accumulation. Independently sampled windows can reset pose or action phase, whereas treating overlap as ordinary noisy
output can rewrite already played history. A single prompt conditions the entire 50-frame
sample, so it also cannot be changed during playback while preserving the completed prefix and regenerating only the
remaining future; the system must wait for the next clip or discard and regenerate the joint sample. Real-time serving
is consequently awkward: playback cannot begin until the joint sample is sufficiently denoised, generation cannot be
naturally pipelined behind playback in bounded blocks, and a prompt update invalidates the prefetched future. Finally,
fixed 50-frame evaluation cannot reveal drift, repetition, or recovery after later continuation boundaries. This is a
task mismatch, not evidence that causal attention is universally superior. Our long-horizon track will instead test a
variable-history block-autoregressive model: completed history is immutable, the prompt may change at each new block,
and frames inside that block still interact bidirectionally, following the windowed distinction used by interactive
motion and chunked-video systems~\cite{zhao2026ardy,magi2025,yin2026evoke}.

\section{Setup}
The testbed, frozen $64^2$ 50-frame protocol, structural oracle (TVR/LIE), and motion diagnostics are those of the
Dancing Stick Figures release~\cite{cho2026dsf}; all real-reference values below come from its motion-disjoint
reference sets. % TODO: compress to ~1 column-inch; cite, don't restate.

% --- moved from testbed report: frozen causal codec paragraph ---
To make block-autoregressive training and 100-frame rollouts affordable we leave pixel space: a 3D causal video VAE
(spatial stride 8, temporal stride 4, 16 latent channels) is trained for 40k steps on the released RGBA frames and
then frozen, with checkpoint and latent statistics pinned by SHA-256. On held-out clips it reconstructs at alpha IoU
$.952$ and $\approx$34\,dB PSNR, and a sliding-window decode audit over 100-frame clips shows no causal-prefix error
or boundary flicker. The higher temporal compression was adopted over a conservative alternative (stride 2, 32
channels) through a pre-declared reconstruction gate rather than throughput preference.

% --- moved from testbed report: M6 architecture description (first half of the M6 paragraph) ---
M6 is a 39.8M full-spatiotemporal latent DiT with 3D rotary embeddings: an immutable clean-history prefix of up to
five latent blocks ($\approx$1\,s) conditions rectified-flow denoising of the next two-latent block (eight video
frames), teacher-forced in training and rolled out to 100 frames at 20\,fps.

\section{Next-block divergence: an AR-native evaluation}
The metric asks the autoregressive question directly: given the data's own history, how far is the model's next block
from the data's actual next block? For each held-out test clip we teacher-force the ground-truth history, generate the
next 16-frame block $N{=}4$ times, and report the best-of-$N$ foreground-weighted RGBA distance to the clip's true
continuation, per block position (32 clips, 6 positions, 10 sampling steps, CFG 2, fixed seed). A \emph{free-running}
variant carries the model's own generated prefix instead, so the teacher-forced-versus-free-running gap measures
exposure bias directly. A \emph{real floor} — the same distance between same-prompt, different-take real clips —
calibrates the scale: it is the divergence of a legitimately different performance. All models score well below this
floor, confirming the metric measures continuation fidelity rather than generic similarity.
Formally, for history $h_b$ ending at block position $b$ and the clip's true continuation $x_b$,
\[
D_{\mathrm{TF}}(b) \;=\; \mathrm{E}_{\mathrm{clips}}\Big[\min_{n\le N} d_{\mathrm{fg}}\big(G(h_b, c, z_n),\, x_b\big)\Big],
\]
where $G$ generates one block under prompt $c$ and noise $z_n$, and $d_{\mathrm{fg}}$ is the union-foreground-weighted
RGBA MSE. $D_{\mathrm{FR}}$ replaces $h_b$ with the model's own rollout prefix. Table~\ref{tab:divergence} reports
both for the pixel-only and rig-co-generating models.

\begin{table}[t]\centering\caption{Teacher-forced next-block divergence (best-of-4) by block position; free-running
(FR) and same-prompt real-pair floor averaged over positions. Lower is better. Matched parameters, codec, data,
optimizer, and recipe; the rig row differs only by co-generated rig tokens.}\label{tab:divergence}\scriptsize
\setlength{\tabcolsep}{3.2pt}
\begin{tabular}{@{}lcccccc c c@{}}
\toprule
model & $b{=}0$ & 1 & 2 & 3 & 4 & 5 & TF avg & FR avg \\
\midrule
pixel-only, 10k & .393 & .369 & .368 & .337 & .232 & .105 & .301 & .465 \\
+rig tokens, 10k & .389 & .361 & .360 & .330 & .227 & .113 & \textbf{.297} & \textbf{.450} \\
pixel-only, 100k & .389 & .370 & .372 & .332 & .251 & .170 & .314 & .458 \\
+rig tokens, 100k & .388 & .360 & .355 & .315 & .213 & .091 & \textbf{.287} & \textbf{.440} \\
\midrule
real-pair floor & .561 & .560 & .563 & .562 & .550 & .531 & .555 & — \\
\bottomrule
\end{tabular}
\end{table}
% NOTE: 10k rig row = canonical-caption v9.0; 100k rig row = varied-caption convergence run (captions are a null
% treatment by the declared caption study, so the pairing is legitimate; state this in the final text).


\section{Co-generating the rig improves next-block prediction}
\todoblock{HEADLINE SECTION. Evidence already collected (results/divergence/, EXPERIMENT\_LOG 2026-08-24): at a
matched 10k budget the rig-co-generating model improves teacher-forced best-of-4 divergence over the pixel-only v8
recipe (.2947 vs .3007) and free-running divergence (.4471 vs .4653); at 100k the gap widens (.2870 vs .3139
teacher-forced, .4399 vs .4584 free-running). Both models share parameters (39.8M), codec, data, optimizer, and the
v8 loss recipe; the ONLY change is appending one rig token per latent frame, denoised jointly. Also striking:
training the pixel-only model $10\times$ longer makes its teacher-forced divergence WORSE (.3007 to .3139) while the
rig-co-generating model improves (.2947 to .2870) — structure tokens change what longer training buys. To write:
seed replication (pending, shares runs with the flagship), rig-quality-vs-pixel-gain analysis, per-block-position
curves, comparison table.}

\section{How tightly should the rig couple to the pixels?}
Co-generation couples rig and pixels only through shared attention. A differentiable soft-capsule renderer can couple
them explicitly: render the predicted clean rig, and match it (foreground-weighted, in premultiplied RGB so limb
identity is visible) against both a soft-render of the ground-truth rig and the decoded predicted pixels, with
gradient flowing into both streams. At constant weight this coupling is a strictly monotone trade: binding rises and
pixel structure falls together (Table~\ref{tab:coupling}), with no weight satisfying both pre-declared gates
(on-figure $\ge$ .85, TVR $\le$ the uncoupled control's .538). A 5k-step linear warmup of the full weight breaks the
trade: applied after the pixel representation settles, full-strength coupling keeps essentially all of the binding
gain (.912 versus the constant arm's .93) while \emph{improving} structure over the uncoupled control (TVR .493
versus .538). Mechanistically this matches the constant arm's failure mode: early in training the decoded pixel
$x_0$ is noise, so the consistency loss drags the pixel path toward an equally noisy rendered rig; ramping the weight
lets attention-level co-generation establish the representation first. All verdicts follow the rule declared in the
sweep protocols before any arm ran.

\begin{table}[t]\centering\caption{Coupling-weight sweep (10k steps each, matched to the uncoupled control; gates
declared before results: on-figure $\ge$ .85 and TVR $\le$ .538). Only the warmed-up arm passes both.}
\label{tab:coupling}\scriptsize
\setlength{\tabcolsep}{4pt}
\begin{tabular}{@{}lcccc@{}}
\toprule
arm & on-figure & TVR & TF div. & FR div. \\
\midrule
uncoupled (v9.0) & .54--.79 & .538 & .297 & .450 \\
$w{=}0.25$ & .724 & \textbf{.518} & .296 & .451 \\
$w{=}0.5$ & .799 & .563 & .295 & .452 \\
$w{=}1.0$ & .876 & .597 & .299 & .463 \\
$w{=}2.0$ constant & \textbf{.930} & .721 & .309 & .497 \\
$w{=}2.0$, 5k warmup & .912 & \textbf{.493} & .298 & .452 \\
\bottomrule
\end{tabular}
\end{table}
% TODO: flagship 100k x 2 seeds of the warmup recipe (launched 2026-08-24) -- add convergence + endpoint.

\section{Rig-space evaluation with a learned regressor (SRE)}
The third layer of the evaluation stack scores any model — rig-free or rig-generating — in rig space. SRE is a
deliberately small instrument: a 3.3M-parameter convolutional regressor mapping one premultiplied RGBA frame to the
27 2-D joints, supervised on the testbed's exact rig labels (loss masked to joints inside the frame) and trained in
minutes on the released cache; its design and validation gates were declared before training. On held-out renders it
is accurate well inside the figure's 1.6\,px capsule radius: 0.66\,px mean joint error on the test split (PCK@2px
$.934$, PCK@4px $.976$; val 0.71\,px). The oracle's corruption harness probes its failure modes. An extra rendered
arm displaces predictions of the duplicated chain by $4.2\times$ the real-frame baseline while other joints stay at
$1.35\times$ — errors are large and localized, not a hallucinated clean skeleton. Swapping the two upper-arm colours
displaces the forearms by $7.7\times$ baseline with the spillover confined, per-joint, to the swapped arm chains
(hands, shoulders, thumbs) and non-arm joints unmoved — the regressor reads the dataset's colour-coded limb identity,
not silhouette alone. (Our pre-declared localization rule counted that chain spillover as non-local and formally
failed this probe; the per-joint decomposition above attributes every displaced joint to the corrupted chains.) The
same harness shows the learned instrument complements the rule-based oracle: a full left--right colour swap,
constructed to preserve adjacency so the oracle cannot flag it, moves SRE's affected joints by $12.8\times$
baseline. One declared gate is vacuous on this data — no released frame carries a joint outside $[-0.5,1.5]$, so
far-off-screen stability is untested; frames with edge-clipped joints degrade gracefully (0.96\,px on test).
\todoblock{Remaining: rig-space next-block divergence table (pixel-only vs rig-co-generating, via SRE on generated
blocks); v9 self-consistency (SRE of decoded pixels vs the model's own rig tokens, at flagship milestones);
bone-length stability under free-running rollout.}

\section{Supporting study I: the static-copy shortcut}
\subsection{The phenomenon}
% --- moved from testbed report: horizon dial + milestone curve + 100k convergence (second half of the M6 paragraph) ---
Small pre-runs expose a generation-horizon dial ($n{=}4$, 2k steps): a four-frame target
block nearly freezes (centroid speed $.120$ versus real $.314$), eight frames is balanced ($.294$), and forty frames
overshoots motion ($.713$) with the worst topology (TVR $.797$). A decoded-RGBA auxiliary loss through the frozen
decoder — differentiable end-to-end, at $1.6\times$ step cost — improved limb existence by 11\% but worsened topology,
occupancy, and prompt sensitivity at $n{=}64$, and was rejected as the default objective. We then trained the balanced
eight-frame model for 20k steps and, by a rule declared before results existed, evaluated every late milestone rather
than the final checkpoint alone ($n{=}64$, paired prompts and noise; Figure~\ref{fig:m6}). Frame structure improves
monotonically (TVR $.745{\to}.349$, LIE $.261{\to}.188$ against real floors of $.133$/$.110$) while motion calibration
collapses early and then flatlines: by step 10k centroid speed falls from $.300$ to $.142$, motion fraction from
$.403$ to $.155$, and pose height variance to one sixth of the real value. Per prompt the model is not uniformly
static — locomotion prompts keep moving while gesture prompts complete one action and rest, echoing the
action-then-rest profile of the data — but the aggregate statistics sit far below matched real clips, so the collapse
is genuine. An independent 100k-step run of the same protocol, evaluated under a convergence rule declared before its
results, settles the question of whether the collapse is under-training: frame structure converges through the real
floors (TVR $.097$ and LIE $.075$ at 100k, below the real $.133$/$.110$) while centroid speed stays pinned near
$.11$--$.13$ from step 10k onward — about a third of the real value — and never recovers (Figure~\ref{fig:m6},
dotted). The model learns to draw ever-cleaner figures that barely move; the freeze is structural, not a budget
artifact.

\begin{figure*}[t]\centering
\includegraphics[width=0.96\textwidth]{paper/figs/m6_milestone_tradeoff.pdf}
\caption{Pre-declared milestone evaluation of the block-autoregressive latent model ($n{=}64$ per checkpoint, paired
prompts and noise, values normalized by the real reference). Structure converges toward its occlusion floor while
motion calibration collapses by step 10k and does not recover; later checkpoints trade motion for anatomy. Dotted
squares show an independent 100k run of the same protocol: structure passes below the real floor while motion stays
collapsed. Diamonds mark the matched full-clip control R0 at its 10k endpoint.}
\label{fig:m6}\end{figure*}

\subsection{Competing hypotheses and matched pilots}
% --- moved from testbed report: "Breaking the freeze" + pilots table ---
Guided by the collapse diagnosis, we ran five matched single-variable 2k-step pilots
against the eight-frame control, each declared before its result (Table~\ref{tab:pilots}). Both history-side
interventions failed — noising the teacher-forced history and extending it to $2.4$\,s both suppress motion further —
while all three interventions that reallocate training signal toward moving figure content revived it:
sixteen-frame generation blocks, a motion-weighted flow loss, and a foreground-alpha loss weight
that upweights latent cells the figure touches (matching the real $.314$ centroid speed at control-level topology; the
transparent background dominates the frame, so uniform latent MSE spends most of its gradient on cells that never
move). The split is mechanistically coherent: the collapse is not a failure to understand history but a gradient
allocation that makes copying rest profitable, so only treatments that change the allocation help.

\begin{table}[t]\centering\caption{Matched single-variable fix pilots (2k steps, $n{=}64$; control: TVR $.745$,
centroid speed $.300$, motion fraction $.403$; real floors $.133$/$.314$/$.371$). Verdicts were declared per
pre-registered config before results existed. Pilot A passes with motion overshoot.}\label{tab:pilots}\scriptsize
\setlength{\tabcolsep}{3pt}
\begin{tabular}{@{}lp{.42\columnwidth}cccc@{}}
\toprule
pilot & treatment & speed & m.\,frac & TVR & verdict \\
\midrule
A & 16-frame generation blocks & .507 & .682 & .760 & pass \\
B & history-noise augmentation $0.2$ & .213 & .276 & — & fail \\
D & motion-weighted flow loss $\alpha{=}1$ & .383 & .523 & .738 & pass \\
E & $2.4$\,s teacher-forced history & .187 & .232 & .748 & fail \\
F & foreground-alpha loss weight $4$ & .327 & .455 & .757 & pass \\
\bottomrule
\end{tabular}
\end{table}

\subsection{The combined repair, replication, and over-training bound}
% --- moved from testbed report: combined fix paragraph ---
A combined run of the three winners at the same 100k budget as the baseline escapes the trade-off
(Figure~\ref{fig:m6fix}): at 100k steps it reaches TVR $.122$ and LIE $.073$ — at the real floors — while holding
centroid speed at $.286$ (91\% of real) and motion fraction at $.434$, where the baseline froze at a third of the
real value. Residual gaps remain in pose height variance ($.028$ versus $.061$) and angular jerk ($1.6\times$ real).
Component attribution rests on the matched pilots. A seed-1 replication of the full 100k run (prediction declared
before results) reproduces the endpoint to within $\pm.005$ on every reported metric (TVR $.122$, speed $.282$), so
the escape is not seed luck. A further $3\times$-budget run
(300k steps, fresh cosine, dotted in Figure~\ref{fig:m6fix}) bounds what training time alone can buy: angular jerk
nearly closes ($.069$ at 300k versus real $.059$) and structure holds below the floor, but centroid speed bleeds
back ($.283$ at 100k to $.245$ at 300k) and the per-prompt motion correlation collapses toward zero while
prompt-versus-noise sensitivity keeps rising — over-training slowly re-freezes the model into smooth, generic,
prompt-agnostic motion. The practical operating range for this recipe is 40k--100k steps; the remaining gaps are
levers for data and conditioning, not for longer schedules.

\begin{figure*}[t]\centering
\includegraphics[width=0.96\textwidth]{paper/figs/m6_fix_comparison.pdf}
\caption{Matched 100k-step runs of the eight-frame baseline and the combined fix (sixteen-frame blocks,
motion-weighted and foreground-weighted flow loss), values normalized by the real reference. Both converge to the
structural floor (a); only the combined fix retains real-level motion (b).}
\label{fig:m6fix}\end{figure*}

\subsection{Attribution: matched full-clip controls}
% --- moved from testbed report: R0 paragraph ---
R0 keeps the frozen codec, backbone family, text encoder, flow objective,
optimizer, and seed, and denoises the complete 25-latent (100-frame) clip jointly with no history prefix for a
pre-declared 10k-step budget. Under the same evaluation it does not freeze: centroid speed is $.382$ and motion
fraction $.540$ — mildly overshooting the real $.314$/$.371$ — with height variance at $.043$ versus M6's $.010$,
while topology is worse than the 20k block-AR model ($.456$ versus $.349$ TVR) and angular jerk is high ($.220$
versus $.059$ real). Since parameters, objective, codec, and data are shared, the motion collapse is attributable to
the teacher-forced short-horizon block factorization rather than to model capacity or the latent flow loss itself,
consistent with the horizon dial above: the shorter the denoised block relative to its clean history, the stronger the
static shortcut. The comparison is diagnostic, not a ranking — R0 trades the collapse for jitter and overshoot, both
runs are single-seed, and R0's budget follows its own declared endpoint. An independent 50k full-clip run under the
same convergence rule confirms the profile at scale: motion stays calibrated throughout (centroid speed
$.375{\to}.293$) but structure plateaus at TVR $\approx.22$ from 30k onward, well short of the $.133$ floor — so
under these budgets the fixed block-autoregressive model above is the only configuration that reaches the structural
floor while keeping real-level motion.

\subsection{Generality: pixel space without a learned codec}
\todoblock{$32^2$ pixel-space no-VAE block-AR pilot. Plan: same block-AR factorization and budget class directly on
$32^2$ RGBA pixels (no learned latent space). Prediction to declare before results: the freeze reproduces (static-copy
is a property of the factorization, not the frozen codec), and the signal-reallocation repair transfers. This section
defuses the ``artifact of the frozen latent space'' objection. Protocol config to be declared before the run.}

\subsection{Residual failure taxonomy}
% --- moved from testbed report: "What remains broken" paragraph ---
Per-prompt analysis of the repaired model splits its
residual errors systematically. Undershoot concentrates on compound vertical motions — squats ($-.47$ centroid-speed
error against matched real clips), burpees ($-.41$), lunges ($-.37$) — which localizes the height-variance deficit;
these numbers are flat from 60k to 100k, so the gap is not a budget problem. Overshoot concentrates on hold-still
prompts (balancing on one leg $+.23$, yoga poses ${\sim}{+}.16$): the model regresses toward mean motion instead of
obeying "stay still". Semantic following saturates: the correlation between per-prompt generated and real motion
statistics rises to ${\sim}.30$ by 40k and then plateaus, while prompt-versus-noise output sensitivity keeps rising
($.46{\to}.79$) — with more training the model listens to the prompt more but does not understand it better. Two
sampler-side diagnostics bound the remaining knobs: raising CFG from 2 to 4 lifts the per-prompt speed correlation
monotonically ($.24{\to}.30$) at no structural cost (TVR improves to ${\sim}.105$), evidence that the conditioning
signal survives the frozen latent space but is weak at the source; and doubling sampler steps to 20 removes about
$40\%$ of the excess angular jerk ($.095{\to}.079$, real $.059$) at mild motion damping, so jerk is partly a sampler
artifact and partly trained-in. A matched pair differing only in the frozen text encoder (t5-small versus t5-base at
10k steps) shows no alignment gain from encoder scale, pointing at the templated one-sentence-per-action captions —
and at conditioning architecture — rather than encoder capacity as the binding constraint. All values are $n{=}64$
under the declared evaluation; per-item details are in the experiment log.

\section{Related work}
\todoblock{Write around: UnityVideo (CVPR 2026, joint multi-modal co-training/co-generation at 10B scale — our scale precedent; full-clip, no continuation measurement, no controlled attribution), LDM3D/IDOL/Structural Video Diffusion (joint RGB+depth/normal), AnimaX (joint video-pose for 3D animation transfer), REPA/V-Co (auxiliary representation alignment/co-denoising), Diffusion Forcing~\cite{chen2024diffusionforcing}, Self Forcing~\cite{huang2025selfforcing}
(exposure bias in AR video diffusion), MAGI-1~\cite{magi2025} (chunk-AR at scale), exposure bias / DAgger literature,
and loss-weighting precedents (min-SNR, P2). Position: those works target the train--test gap at rollout time; the
static-copy shortcut is a \emph{training-time} gradient-allocation failure that appears before any rollout drift, and
history-side remedies do not touch it. The frontier-state paragraph below is drafted; fold it in rather than
re-covering EVOKE.}

% --- EVOKE positioning: frontier evidence for explicit structural state in chunk-AR generation ---
\paragraph{Explicit structural state at the frontier.} Concurrent interactive systems arrive at this paper's design
axis from the deployment side. Alaya-EVOKE~\cite{yin2026evoke} — a 14B chunk-autoregressive diffusion world model
distilled to three CFG-free steps — keeps scene geometry in a bounded, camera-indexed state bank outside the
denoiser, estimated by a monocular depth model and re-rendered into each chunk's conditioning, precisely so that
hours-long rollouts need not trust the pixel stream to carry structure. The premise is the one this paper tests:
a block-autoregressive video generator drifts less when explicit structural state travels alongside the pixels. The
mechanisms are complementary — EVOKE \emph{conditions on estimated} state where our model \emph{co-generates exact}
state — and so is the evidence: a deployed system's benchmark scores (WBench, VBench-Long) aggregate heuristic
quality with no matched state-free control at scale, and cannot ask whether a given continuation matches the world
it extends. Our testbed supplies the causal form of the question such systems build on: with the rig exact and
parameters, data, and budget matched, the structural channel by itself improves next-block prediction — measured
against the data's actual continuation.

\section{Limitations}
\todoblock{Single traced synthetic domain (defended, not hidden); single-seed pilots with seed-replicated main run;
one backbone family and codec; metric suite is image-space diagnostics; alignment remains open (taxonomy section) —
enumerate honestly, mirroring the testbed report's style.}

\section{Conclusion}
\todoblock{Write last.}

\paragraph{AI assistance disclosure.} Pixel (\href{https://sprited.ai}{@pix-el}) and Claudia assisted; the author
reviewed the work and is responsible for its claims.

\balance
\begin{thebibliography}{99}\small
\bibitem{cho2026dsf} J. H. Cho. Dancing Stick Figures: A Traced Synthetic Testbed for Video Generation
Diagnostics. arXiv, 2026. % TODO: fill arXiv id after the testbed report is submitted.
\bibitem{magi2025} Sand AI et al. MAGI-1: Autoregressive Video Generation at Scale. arXiv:2505.13211, 2025.
\bibitem{chen2024diffusionforcing} B. Chen et al. Diffusion Forcing: Next-token Prediction Meets Full-Sequence Diffusion. \emph{NeurIPS}, 2024.
\bibitem{huang2025selfforcing} X. Huang et al. Self Forcing: Bridging the Train-Test Gap in Autoregressive Video Diffusion. arXiv:2506.08009, 2025.
\bibitem{zhao2026ardy} K. Zhao et al. ARDY: Autoregressive Diffusion with Hybrid Representation for Interactive Human Motion Generation. arXiv:2607.08741, 2026.
\bibitem{yin2026evoke} Y. Yin et al. Alaya-EVOKE: From Linear-Scaling Supervision to Endless World. arXiv:2608.13546, 2026.
\end{thebibliography}
\end{document}
